\question[12] به سوالات زیر پاسخ مناسب دهید.

\begin{parts}

    \part ساختار داده مناسب برای پیاده‌سازی بافر محدود
    \lr{(Bounded Buffer)}
    چیست؟

    \part دو فرآیند برای برقراری ارتباط از طریق حافظه اشتراکی، باید چه محدودیت خاصی که توسط سیستم‌عامل اعمال می‌شود را حذف کنند؟

    \part فرآیندها چه تضمینی را باید در مورد نوشتن در مکان حافظه اشتراکی ارائه دهند؟

    \part    پیوند ارتباطی به چه روش‌هایی می‌تواند به صورت فیزیکی پیاده‌سازی شود؟

    \part هنگامی که مالکیت صندوق پستی با سیستم عامل است، سیستم عامل چه سه مکانیزمی را باید فراهم کند؟

    \part مزیت «پاسخگویی سریع» در برنامه‌های چندنخی به چه معناست؟

    \part یک نخ (Thread) به عنوان واحد پایه استفاده از
    \lr{CPU}،
    شامل چه ۴ جزء اختصاصی است؟

    \part چه منابعی بین تمام نخ‌های یک فرآیند مشترک هستند؟ (۳ مورد نام ببرید)

    \part چهار مزیت اصلی برنامه‌نویسی چندنخی که در اسلایدها ذکر شده است را نام ببرید.

    \part پنج چالش اصلی برنامه‌نویسی در سیستم‌های چند‌هسته‌ای را نام ببرید.

    \part چرا مدل «یک به یک» هم‌روندی (Concurrency) بیشتری نسبت به مدل «چند به یک» دارد؟

    \part دو مشکل اصلی ایجاد نخ جدید برای هر درخواست (بدون استفاده از استخر نخ) را نام ببرید.

    \part تفاوت اصلی TLS با داده‌های ایستا
    \lr{(Static Data)}
    چیست؟
\end{parts}